\documentclass[12pt, letterpaper]{article}

% --- PAQUETES ---
\usepackage[utf8]{inputenc}
\usepackage[T1]{fontenc}
\usepackage[spanish]{babel}
\usepackage[left=2cm, right=2cm, top=2cm, bottom=2.5cm]{geometry}
\usepackage{enumitem}
\usepackage{amssymb}
\usepackage{tcolorbox}
\usepackage{amsmath}

% --- FUENTE --- 
\usepackage{helvet} 
\renewcommand{\familydefault}{\sfdefault}

\begin{document}
	
	% --- ENCABEZADO --- 
	\begin{center} 
		\textbf{Taller de Repaso - CIENCIAS 2}\\[0.1cm] 
		\textit{TAU 1: ¡Viajemos por dentro de la Tierra!} 
	\end{center}
	
	
	% --- PREGUNTA 1 --- 
	\noindent \textbf{1. Viajeros de la mente:} El libro menciona que la segunda mejor forma de viajar por el Universo (después de una nave espacial) es usando la mente y la ciencia. Nombra a los dos famosos científicos mencionados en el TAU que lograron grandes descubrimientos sobre el Universo de esta manera:
	\vspace{0.2cm}
	\begin{itemize}[itemsep=0.4cm, label={--}]
		\item \underline{\hspace{8cm}}
		\item \underline{\hspace{8cm}}
	\end{itemize}
	
	\vspace{0.3cm}
	
	% --- PREGUNTA 2 --- 
	\noindent \textbf{2. Saliendo de la Tierra:} Cuando el transbordador espacial sube y se aleja a unos 2.500 kilómetros de altura, la ilusión óptica de ver la Tierra como un ``plato de sopa'' desaparece. Escribe brevemente cómo se ve nuestro planeta realmente desde esa altura en el espacio:
	\vspace{0.8cm}
	\noindent \rule{\textwidth}{0.4pt}
	\vspace{0.4cm}
	
	% --- PREGUNTA 3 --- 
	\noindent \textbf{3. Dimensiones y cálculos de nuestro planeta:} Completa los datos en blanco sobre el tamaño de la Tierra y resuelve el problema matemático.
	\vspace{0.2cm}
	\begin{itemize}[itemsep=0.4cm, label={$\bullet$}]
		\item La distancia desde la superficie hasta el \textbf{centro exacto} de la Tierra es de \underline{\hspace{2.5cm}} kilómetros.
		\item La circunferencia (también llamada ``cintura'') de la Tierra mide \underline{\hspace{2.5cm}} kilómetros.
	\end{itemize}
	\vspace{0.2cm}
	\noindent \textbf{Problema:} Si decidieras darle la vuelta a la Tierra \textbf{caminando} y te mueves a 5 $km/h$ durante 10 $h$ cada día, ¿cuántos días te demorarías en dar la vuelta completa a la tierra?
	\vspace{0.1cm}

	\vspace{0.2cm}
	\begin{tcolorbox}[colframe=black, colback=white, height=3cm, width=\textwidth]
	\end{tcolorbox}
	\vspace{0.4cm}
	\newpage
	% --- PREGUNTA 4 --- 
	\noindent \textbf{4. La Tierra por dentro:} La geofísica nos ha enseñado que el interior del planeta tiene zonas sólidas y líquidas. Relaciona con una línea cada capa de la Tierra con su descripción correcta:
	
	\vspace{0.4cm}
	\begin{itemize}[itemsep=0.8cm, label={}]
		\item \textbf{Corteza} \hfill $\circ$ Capa que llega hasta el centro, formada por materia sólida.
		\item \textbf{Manto} \hfill $\circ$ Es la "piel" sólida de la Tierra, donde vivimos y están los océanos.
		\item \textbf{Núcleo exterior} \hfill $\circ$ Capa formada por materia totalmente \textbf{líquida}, caliente y pesada.
		\item \textbf{Núcleo interior} \hfill $\circ$ Capa de roca sólida y caliente que va de los 40 a los 3.000 km.
	\end{itemize}
	
	\vspace{0.5cm}
	
	% --- PREGUNTA 5 --- 
	\noindent \textbf{5. Una historia abisal:} En la última sección del libro viajamos a lo más profundo del océano. Responde Falso (F) o Verdadero (V) a las siguientes afirmaciones:
	
	\vspace{0.3cm}
	\begin{enumerate}[label=\alph*), itemsep=0.4cm]
		\item La palabra ``abisal'' viene del griego \textit{ábyssos} y signfica ``sin fondo''. \hfill ( \hspace{0.5cm} )
		\item El pequeño y súper-reforzado submarino que bajó a la Fosa de las Marianas en 1960 se llamaba \textit{Trieste}. \hfill ( \hspace{0.5cm} )
		\item En el fondo de la Fosa de las Marianas no se encontró ningún ser vivo, porque la presión del agua aplasta cualquier cuerpo vertebrado. \hfill ( \hspace{0.5cm} )
	\end{enumerate}
	
	
\end{document}
```

**Guía rápida para la calificación del taller:**
*   **Punto 1:** Albert Einstein y Stephen Hawking.
*   **Punto 2:** Como una hermosa esfera brillante y blanquiazul, flotando en el vacío.
*   **Punto 3:** Centro: 6.400 km. Cintura: 40.000 km. Cálculo: $40.000 \div 50 = 800$ días.
*   **Punto 4:** Corteza (piel sólida) / Manto (roca sólida hasta 3000 km) / Núcleo exterior (materia líquida) / Núcleo interior (llega hasta el centro, materia sólida).
*   **Punto 5:** a) Verdadero / b) Verdadero / c) Falso (sí encontraron vida: un pez plano).